\documentclass[11pt]{article}
\newcommand{\ListaTitulo}{Lista 08: Simulado Prova 1 (Unidade 1)}
% Preâmbulo compartilhado: MATF14, Listas de Exercícios 2026
% Usado por ListaNN.tex e GabaritoNN.tex via % Preâmbulo compartilhado: MATF14, Listas de Exercícios 2026
% Usado por ListaNN.tex e GabaritoNN.tex via % Preâmbulo compartilhado: MATF14, Listas de Exercícios 2026
% Usado por ListaNN.tex e GabaritoNN.tex via \input{preamble}

\usepackage[utf8]{inputenc}
\usepackage[T1]{fontenc}
\usepackage[brazil]{babel}
\usepackage{fullpage}
\usepackage{amsmath,amssymb}
\usepackage{enumitem}
\usepackage{xcolor}
\usepackage{fancyhdr}
\usepackage{titlesec}
\usepackage{listings}
\usepackage{hyperref}
\usepackage{booktabs}
\usepackage{graphicx}
\usepackage{tikz}

\definecolor{matf14azul}{HTML}{0D3B54}
\definecolor{matf14dourado}{HTML}{C9971E}

\titleformat{\section}{\normalfont\Large\bfseries\color{matf14azul}}{\thesection}{1em}{}
\titleformat{\subsection}{\normalfont\large\bfseries\color{matf14azul}}{\thesubsection}{1em}{}

\providecommand{\ListaTitulo}{Lista de Exercícios}

\setlength{\headheight}{14pt}
\pagestyle{fancy}
\fancyhf{}
\lhead{\small MATF14: Estatística Econômica I}
\rhead{\small \ListaTitulo}
\cfoot{\thepage}
\renewcommand{\headrulewidth}{0.4pt}

\lstset{
  basicstyle=\ttfamily\small,
  breaklines=true,
  frame=single,
  columns=fullflexible,
  backgroundcolor=\color{gray!8},
  commentstyle=\color{matf14dourado},
  keywordstyle=\color{matf14azul}\bfseries,
  language=R
}

\newcounter{questaocounter}
\newenvironment{questao}[1]{%
  \stepcounter{questaocounter}%
  \bigskip\noindent\textbf{Questão #1.}\ %
}{\par}

\newcommand{\cabecalho}[2]{%
  \begin{center}
    {\Large\bfseries\color{matf14azul} #1}\\[0.3em]
    {\large #2}
  \end{center}
  \vspace{0.5em}
  \noindent\rule{\linewidth}{0.6pt}
  \vspace{0.5em}
}

\newenvironment{solucao}{%
  \par\smallskip\noindent\textit{\color{matf14dourado!80!black}Solução.}\ %
}{\hfill$\blacksquare$\par}


\usepackage[utf8]{inputenc}
\usepackage[T1]{fontenc}
\usepackage[brazil]{babel}
\usepackage{fullpage}
\usepackage{amsmath,amssymb}
\usepackage{enumitem}
\usepackage{xcolor}
\usepackage{fancyhdr}
\usepackage{titlesec}
\usepackage{listings}
\usepackage{hyperref}
\usepackage{booktabs}
\usepackage{graphicx}
\usepackage{tikz}

\definecolor{matf14azul}{HTML}{0D3B54}
\definecolor{matf14dourado}{HTML}{C9971E}

\titleformat{\section}{\normalfont\Large\bfseries\color{matf14azul}}{\thesection}{1em}{}
\titleformat{\subsection}{\normalfont\large\bfseries\color{matf14azul}}{\thesubsection}{1em}{}

\providecommand{\ListaTitulo}{Lista de Exercícios}

\setlength{\headheight}{14pt}
\pagestyle{fancy}
\fancyhf{}
\lhead{\small MATF14: Estatística Econômica I}
\rhead{\small \ListaTitulo}
\cfoot{\thepage}
\renewcommand{\headrulewidth}{0.4pt}

\lstset{
  basicstyle=\ttfamily\small,
  breaklines=true,
  frame=single,
  columns=fullflexible,
  backgroundcolor=\color{gray!8},
  commentstyle=\color{matf14dourado},
  keywordstyle=\color{matf14azul}\bfseries,
  language=R
}

\newcounter{questaocounter}
\newenvironment{questao}[1]{%
  \stepcounter{questaocounter}%
  \bigskip\noindent\textbf{Questão #1.}\ %
}{\par}

\newcommand{\cabecalho}[2]{%
  \begin{center}
    {\Large\bfseries\color{matf14azul} #1}\\[0.3em]
    {\large #2}
  \end{center}
  \vspace{0.5em}
  \noindent\rule{\linewidth}{0.6pt}
  \vspace{0.5em}
}

\newenvironment{solucao}{%
  \par\smallskip\noindent\textit{\color{matf14dourado!80!black}Solução.}\ %
}{\hfill$\blacksquare$\par}


\usepackage[utf8]{inputenc}
\usepackage[T1]{fontenc}
\usepackage[brazil]{babel}
\usepackage{fullpage}
\usepackage{amsmath,amssymb}
\usepackage{enumitem}
\usepackage{xcolor}
\usepackage{fancyhdr}
\usepackage{titlesec}
\usepackage{listings}
\usepackage{hyperref}
\usepackage{booktabs}
\usepackage{graphicx}
\usepackage{tikz}

\definecolor{matf14azul}{HTML}{0D3B54}
\definecolor{matf14dourado}{HTML}{C9971E}

\titleformat{\section}{\normalfont\Large\bfseries\color{matf14azul}}{\thesection}{1em}{}
\titleformat{\subsection}{\normalfont\large\bfseries\color{matf14azul}}{\thesubsection}{1em}{}

\providecommand{\ListaTitulo}{Lista de Exercícios}

\setlength{\headheight}{14pt}
\pagestyle{fancy}
\fancyhf{}
\lhead{\small MATF14: Estatística Econômica I}
\rhead{\small \ListaTitulo}
\cfoot{\thepage}
\renewcommand{\headrulewidth}{0.4pt}

\lstset{
  basicstyle=\ttfamily\small,
  breaklines=true,
  frame=single,
  columns=fullflexible,
  backgroundcolor=\color{gray!8},
  commentstyle=\color{matf14dourado},
  keywordstyle=\color{matf14azul}\bfseries,
  language=R
}

\newcounter{questaocounter}
\newenvironment{questao}[1]{%
  \stepcounter{questaocounter}%
  \bigskip\noindent\textbf{Questão #1.}\ %
}{\par}

\newcommand{\cabecalho}[2]{%
  \begin{center}
    {\Large\bfseries\color{matf14azul} #1}\\[0.3em]
    {\large #2}
  \end{center}
  \vspace{0.5em}
  \noindent\rule{\linewidth}{0.6pt}
  \vspace{0.5em}
}

\newenvironment{solucao}{%
  \par\smallskip\noindent\textit{\color{matf14dourado!80!black}Solução.}\ %
}{\hfill$\blacksquare$\par}


\begin{document}

\cabecalho{Lista de Exercícios 08: Simulado}{Revisão geral, no formato e no nível da Prova 1: Estatística Descritiva (Cap. 1, Aulas 01--08)}

\noindent \textit{Simulado de 10 pontos, no mesmo formato da prova oficial: questões com múltiplos
itens, valor indicado entre parênteses. Resolva com tempo cronometrado (sugestão: 100 minutos,
sem consulta) antes de olhar o gabarito em \texttt{Gabarito08.pdf}.}

\begin{enumerate}

\item Uma amostra de 40 pequenas empresas de um município foi classificada por setor de atividade:
28 em Comércio, 8 em Indústria e 4 em Serviços. Considere também que, das 40 empresas, o porte
declarado foi: 22 Pequena, 14 Média e 4 Grande.

\begin{enumerate}
  \item (1 ponto) Classifique as variáveis "setor de atividade" e "porte" (nominal, ordinal,
    discreta ou contínua).
  \item (1 ponto) Construa a tabela de frequências absolutas e relativas (\%) para a variável
    setor de atividade.
  \item (1 ponto) Qual tipo de gráfico você recomendaria para apresentar a distribuição por setor a
    uma associação comercial? Justifique frente a pelo menos uma alternativa descartada.
\end{enumerate}

\item A tabela abaixo traz o número de reclamações mensais recebidas por 12 lojas de uma rede de
varejo:

\begin{center}
\begin{tabular}{cccccccccccc}
\toprule
3 & 5 & 4 & 20 & 6 & 4 & 5 & 7 & 3 & 6 & 5 & 4 \\
\bottomrule
\end{tabular}
\end{center}

\begin{enumerate}
  \item (1,5 ponto) Calcule a média, a mediana e a moda do número de reclamações.
  \item (1 ponto) Calcule o desvio padrão (divisor $n$).
  \item (1,5 ponto) Uma das lojas (com 20 reclamações) está passando por uma crise pontual de
    atendimento. Interprete o efeito dessa loja sobre cada uma das três medidas de tendência
    central calculadas em (a), e indique qual delas você recomendaria usar para descrever "a loja
    típica da rede".
\end{enumerate}

\item Uma corretora registrou, em 8 meses, o gasto mensal com marketing digital ($X$, em R\$ mil)
e o número de novos clientes captados no mês seguinte ($Y$):

\begin{center}
\begin{tabular}{c|cccccccc}
\toprule
$X$ & 4 & 6 & 5 & 8 & 7 & 9 & 5 & 10 \\
\midrule
$Y$ & 12 & 18 & 15 & 22 & 20 & 25 & 14 & 27 \\
\bottomrule
\end{tabular}
\end{center}

\begin{enumerate}
  \item (1,5 ponto) Calcule o coeficiente de correlação de Pearson entre $X$ e $Y$ e interprete o
    resultado em uma frase.
  \item (1 ponto) O gerente de marketing conclui que "gastar mais em marketing causa mais clientes
    novos". Aponte uma explicação alternativa (não causal) plausível para essa correlação.
\end{enumerate}

\item O faturamento mensal (em R\$ mil) de uma amostra de 50 microempresas tem média R\$ 42 mil e
mediana R\$ 31 mil.

\begin{enumerate}
  \item (0,5 ponto) A distribuição do faturamento é aproximadamente simétrica, assimétrica à
    direita ou assimétrica à esquerda? Justifique com base na relação entre média e mediana.
  \item (1 ponto) Um consultor quer comunicar "o faturamento típico" dessas microempresas para um
    público leigo. Qual medida (média ou mediana) você recomendaria, e por quê?
\end{enumerate}

\end{enumerate}

\vspace{1cm}
\begin{center}
\Large\textbf{Formulário}
\end{center}
\begin{small}
\begin{itemize}
  \item $\bar x = \frac1n\sum_i x_i$; mediana $=x_{((n+1)/2)}$ ($n$ ímpar) ou
    $\frac{x_{(n/2)}+x_{(n/2+1)}}{2}$ ($n$ par).
  \item $\mathrm{Var}(X)=\frac1n\sum_i(x_i-\bar x)^2$; $dp(X)=\sqrt{\mathrm{Var}(X)}$;
    $CV=dp(X)/\bar x$.
  \item $r_{XY} = \dfrac{\sum_i (x_i-\bar x)(y_i-\bar y)}{\sqrt{\sum_i(x_i-\bar x)^2}\sqrt{\sum_i(y_i-\bar y)^2}}$.
  \item Assimetria $=\dfrac{\frac1n\sum_i(x_i-\bar x)^3}{\mathrm{Var}(X)^{3/2}}$; moda $<$ mediana
    $<$ média sob assimetria positiva (e vice-versa sob assimetria negativa).
\end{itemize}
\end{small}

\end{document}
